% Official CC BY4.0 published HTML S5; original math alttext and author text, literal original numbering.


\par\medskip\noindent\textbf{5 Efficient construction of projectors}\par




In this section we prove Theorem 1.1.
We explain the construction of the projectors in Section 5.1, and describe the efficient circuit implementation in Section 5.2.





\par\medskip\noindent\textbf{5.1 Construction of projectors}\par




We construct the required projections as projections onto a common eigenspace of several commuting matrices. The construction is elementary and is based on the classical construction of eigenprojectors from basic linear algebra.





We present the construction in full generality in Theorem 5.2. The projectors in the cases relevant for Theorem 1.1 are obtained by applying the general construction to suitable commutative subalgebras of the universal enveloping algebra $\mathcal{U}(\mathfrak{gl}_{k})$, see Table 1. If $\mathcal{A}$ is a commutative algebra generated by $U_{1},\dots,U_{q}$, then a representation of $\rho:\mathcal{A}\to\textnormal{End}(V)$ is uniquely determined by $\rho(U_{1}),\ldots,\rho(U_{q})$, which are $q$ commuting endomorphisms of $V\to V$.
Every irreducible representation of $\mathcal{A}$ is $1$-dimensional, so it is given by an algebra homomorphism $\chi\colon\mathcal{A}\to\mathbb{C}$.
In other words, an irreducible representation is spanned by a common eigenvector of the generators $U_{1},\dots,U_{q}$ of $\mathcal{A}$ with the corresponding eigenvalues $\chi(U_{1}),\dots,\chi(U_{q})$.





Recall that a representation is called completely reducible if it can be decomposed into a direct sum of irreducible representations.
All representations in this paper have this property.
For a commutative algebra this means that the generators of the algebra act not just by commuting but by simultaneously diagonalizable linear maps.





We slightly abuse notation and denote the character $\chi\colon\mathcal{A}\to\mathbb{C}$ and the corresponding irreducible representation by the same letter.
Every completely reducible representation $V$ of $\mathcal{A}$ has a unique isotypic decomposition



\[\begin{aligned}
\displaystyle V\cong\bigoplus_{\lambda\in\Lambda}\chi_{\lambda}^{\oplus m_{\lambda}},\qquad\text{(15)}
\end{aligned}\]


where $\Lambda$ is a set indexing irreducible representations of $\mathcal{A}$.
We denote the isotypic component $\chi_{\lambda}^{\oplus m_{\lambda}}$ in $V$ by $V_{\mathcal{A},\lambda}$.





The length of the projectors as elements of $\mathcal{U}(\mathfrak{g})$ depends mainly on the number of different isotypic components of the representation.





\par\medskip\noindent\textbf{Definition 5.1.}\par



Let $V$ be a completely reducible representation of an algebra $\mathcal{A}$.
We call the number of isomorphism types of irreducible subrepresentations of $V$ the diversity of $V$.






\par\medskip\noindent\textbf{Theorem 5.2.}\par



Let $\mathcal{A}$ be a commutative algebra generated by $U_{1},\dots,U_{q}$ and $V$ be a completely reducible representation of $\mathcal{A}$ with diversity at most $p$.
For every irreducible representation $\chi_{\lambda}$ of $\mathcal{A}$ there exists a polynomial $f_{\lambda}$ of degree at most $p-1$ such that



\[\begin{aligned}
\displaystyle f_{\lambda}(U_{1},\dots,U_{q}).x=\Pi_{\lambda}x\qquad\text{(16)}
\end{aligned}\]


where $\Pi_{\lambda}\colon V\to V$ is the unique $\mathcal{A}$-equivariant projection onto the isotypic component $V_{\mathcal{A},\lambda}$.






\par\medskip\noindent\textbf{Proof.}\par




Let $\Lambda$ be the set of isomorphism types appearing in the isotypic decomposition of $V$, so that $V\cong\bigoplus_{\mu\in\Lambda}V_{\mathcal{A},\mu}$ with $V_{\mathcal{A},\mu}\neq 0$. An element $U\in\mathcal{A}$ acts as $U.x=\chi_{\mu}(U)x$ on $V_{\mathcal{A},\mu}$. If $\lambda\notin\Lambda$, set $f_{\lambda}=0$.





Fix $\lambda\in\Lambda$. For every $\mu\in\Lambda$ with $\mu\neq\lambda$, let $U_{i_{\mu}}$ be a generator such that $\chi_{\mu}(U_{i_{\mu}})\neq\chi_{\lambda}(U_{i_{\mu}})$. Such generator exist because $\chi_{\lambda}$ and $\chi_{\mu}$ are different characters of $\mathcal{A}$. Define



\[\begin{aligned}
\displaystyle P_{\lambda}\coloneqq\prod_{\mu\in\Lambda,\mu\neq\lambda}\frac{U_{i_{\mu}}-\chi_{\mu}(U_{i_{\mu}})}{\chi_{\lambda}(U_{i_{\mu}})-\chi_{\mu}(U_{i_{\mu}})}\quad\in\mathcal{A}.\qquad\text{(17)}
\end{aligned}\]


The factor $\dfrac{U_{i_{\mu}}-\chi_{\mu}(U_{i_{\mu}})}{\chi_{\lambda}(U_{i_{\mu}})-\chi_{\mu}(U_{i_{\mu}})}$ acts on $V$ as the zero map on $V_{\mathcal{A},\mu}$ and as the identity map on $V_{\mathcal{A},\lambda}$; moreover, its action maps isotypic components to themselves. The element $P_{\lambda}$ acts as the composition of these factors; therefore its image in $\textnormal{End}(V)$ is the projection $\Pi_{\lambda}$ onto $V_{\mathcal{A},\lambda}$. Finally, since every factor is linear in $U_{1},\ldots,U_{q}$, $P_{\lambda}$ is a polynomial in $U_{1},\ldots,U_{q}$. Its degree is (at most) $p-1$.
$\hfill\blacktriangleleft$





The projections required in Theorem 1.1 arise as special cases of this construction.
In each case, the image of the projection is an isotypic component of some commutative subalgebra $\mathcal{A}\subset\mathcal{U}(\mathfrak{gl}_{k})$.
Theorem 5.2 then gives a construction of the projectors as elements of $\mathcal{U}(\mathfrak{gl}_{k})$, expressed as polynomials in the generators of $\mathcal{A}$.
The length of the projectors is easily bounded in terms of lengths of generators of $\mathcal{A}$ in $\mathcal{U}(\mathfrak{gl}_{K})$ and the diversity of $\mathbb{C}[\mathbb{C}[x_{1},\ldots,x_{k}]_{d}]_{\delta}$ as a representation of $\mathcal{A}$.
Indeed, if the generators $U_{1},\dots,U_{q}$ of $\mathcal{A}$ have length at most $\ell$, and the projectors are expressed as polynomials in $U_{1},\dots,U_{q}$ of degree at most $p-1$, then the lengths of the projectors are bounded by $\ell(p-1)$. The following table summarizes these parameters for the three cases considered.




Table 1: 
Isotypic components, generators of their corresponding subalgebras, and their diversity in $\mathbb{C}[\mathbb{C}[x_{1},\ldots,x_{k}]_{d}]_{\delta}$.
See Equation 8 and Equation 14 for the definitions of $C_{p}$ and $C_{\ell,p}$.


\begin{quote}
Isotypic components \quad;\quad Weight spaces \quad;\quad $\lambda$-isotypic spaces \quad;\quad $T$-isotypic spaces\par
Algebra $\mathcal{A}$ \quad;\quad $\mathcal{U}(\mathfrak{h})\cong\mathbb{C}[\mathfrak{h}^{*}]$ \quad;\quad $\mathcal{Z}(\mathfrak{gl}_{k})$ \quad;\quad $\mathcal{G}\hskip-1.00006pt\mathcal{Z}(\mathfrak{gl}_{k})$\par
Generators $U_{i}$ \quad;\quad $\big{\{}E_{i,i}\mid i\in[k]\big{\}}$ \quad;\quad $\big{\{}C_{p}\mid p\in[k]\big{\}}$ \quad;\quad $\big{\{}C_{\ell,p}\mid\ell\in[k],p\in[\ell]\big{\}}$\par
Generator lengths \quad;\quad $1$ \quad;\quad $\leq k$ \quad;\quad $\leq k$\par
Diversity \quad;\quad $\leq(\delta d)^{k}$ \quad;\quad $\leq(\delta d)^{k}$ \quad;\quad $\leq(\delta d)^{k^{2}}$\par
\end{quote}




\par\medskip\noindent\textbf{$\blacktriangleright$ Remark 5.3.}\par



In the construction of Theorem 5.2 we require knowledge of the value of $\chi_{\mu}$ evaluated on the generators $U_{i}$, for all irreducible representation types $\mu$ occurring in the representation.
For the weight spaces these values are given by the weights themselves, and for the other two algebra they can be computed using Equation 9.
It is then also possible to remove any choices in the construction by taking for each term instead a linear combination $U$ of all $U_{i}$ with algebraically independent irrational weights.
Since the values $\chi_{\mu}$ are always integer, this ensures $\chi_{\mu}(U)\neq\chi_{\lambda}(U)$ when $\mu\neq\lambda$.






\par\medskip\noindent\textbf{Example 5.4.}\par



We construct and evaluate a $\lambda$-isotypic projector for metapolynomials of format $(3,2,3)$.
From Example 4.8 we know $\Lambda=\{(6,0,0),(4,2,0),(2,2,2)\}$.
By Equation 9 we find that the Casimir element $C_{2}$ scales these three irreducible representations with different scalars, namely $48$, $28$ and $12$ respectively.
Hence we can use $C_{2}$ to construct the projectors from Theorem 5.2.
Consider first the projector $P_{(6,0,0)}$ to the $(6,0,0)$-isotypic component.
We have that



\[\begin{aligned}
\displaystyle P_{(6,0,0)}=\frac{C_{2}-\chi_{(2,2,2)}(C_{2})}{\chi_{(6,0,0)}(C_{2})-\chi_{(2,2,2)}(C_{2})}\cdot\frac{C_{2}-\chi_{(4,2,0)}(C_{2})}{\chi_{(6,0,0)}(C_{2})-\chi_{(4,2,0)}(C_{2})}=\frac{C_{2}-12}{36}\cdot\frac{C_{2}-28}{20}.
\end{aligned}\]


Consider again the example $\Delta=c_{002}c_{020}c_{200}$ from the introduction.
We will compute $P_{(6,0,0)}.\Delta$. Using the equations in ˜4.2 we find



\[\begin{aligned}
\displaystyle C_{2}.\Delta & \displaystyle=\sum_{i=1}^{3}\sum_{j=1}^{3}\big{(}E_{i,j}E_{j,i}\big{)}.\Delta=2\,c_{020}c_{101}^{2}+2\,c_{002}c_{110}^{2}+2\,c_{011}^{2}c_{200}+24\,c_{002}c_{020}c_{200}.
\end{aligned}\]


It directly follows that



\[\begin{aligned}
\displaystyle C_{2}.\Delta-28\Delta=2\,c_{020}c_{101}^{2}+2\,c_{002}c_{110}^{2}+2\,c_{011}^{2}c_{200}-4\,c_{002}c_{020}c_{200}.
\end{aligned}\]


We apply $C_{2}$ again (at this point a computer algebra program is recommended) and find



\[\begin{aligned}
\displaystyle C_{2}.(C_{2}.\Delta-28\Delta)=48\big{(}c_{020}c_{101}^{2}+c_{011}c_{101}c_{110}+c_{002}c_{110}^{2}+c_{011}^{2}c_{200}\big{)}.
\end{aligned}\]


Hence,



\[\begin{aligned}
\displaystyle 60\,P_{(6,0,0)}.\Delta & \displaystyle=\frac{C_{2}.(C_{2}.\Delta-28\Delta)-12(C_{2}.\Delta-28\Delta)}{12} \\
 & \displaystyle=2\,c_{020}c_{101}^{2}+4\,c_{011}c_{101}c_{110}+2\,c_{002}c_{110}^{2}+2\,c_{011}^{2}c_{200}+4\,c_{002}c_{020}c_{200}.
\end{aligned}\]


This is indeed the isotypic component as given in the introduction.
The other two can be determined by an analogous computation.







\par\medskip\noindent\textbf{5.1.1 Weight projectors}\par




Let $V$ be a representation of a complex reductive Lie group $G$.
The action of $G$ on $V$ induces the action of the corresponding Lie algebra $\mathfrak{g}$.
Fix a maximal torus $T\subset G$ and let $\mathfrak{h}\subset\mathfrak{g}$ be the corresponding Cartan subalgebra.
The torus $T$ is reductive, and the isotypic components of $V$ with respect to the action of $T$ are exactly the weight spaces.
The action of $T$ induces the action of $\mathfrak{h}$ with the same isotypic components, so the action of $\mathfrak{h}$ is completely reducible.





Let $H_{1},\dots,H_{r}$ be a basis of the Cartan subalgebra $\mathfrak{h}$.
Since $\mathfrak{h}$ is abelian, the universal enveloping algebra $\mathcal{U}(\mathfrak{h})\subset\mathcal{U}(\mathfrak{g})$ is commutative.
It can be identified with the polynomial algebra $\mathbb{C}[H_{1},\dots,H_{r}]$.
We can now apply Theorem 5.2 to obtain for every weight $\mu\in\mathfrak{h}^{*}$ an element $H_{\mu}\in\mathcal{U}(\mathfrak{g})$ such that $H_{\mu}.v$ is the projection of $v\in V$ onto the weight space of weight $\mu$.
The projectors $H_{\mu}$ are obtained as polynomial expressions $H_{\mu}=f_{\mu}(H_{1},\dots,H_{r})$ with $\deg f_{\mu}\leq p-1$ where $p$ is the diversity of $V$ as a representation of $\mathfrak{h}$ (or, equivalently, of $\mathcal{U}(\mathfrak{h})$), that is, the number of weights of $V$.
Since $H_{i}$ as elements of $\mathcal{U}(\mathfrak{g})$ have length $1$, the length of $H_{\mu}$ is also bounded by $p-1$.





We can specialize this construction to the action of $\mathfrak{gl}_{k}$ on the space of metapolynomials $W\coloneqq\mathbb{C}[\mathbb{C}[x_{1},\ldots,x_{k}]_{d}]_{\delta}$.
The standard Cartan subalgebra consists of all diagonal matrices, and we can take $H_{i}\coloneqq E_{i,i}$.
The weights appearing in $W$ correspond to $k$-tuples of nonnegative integers summing to $\delta d$ (see Example 4.7). The number of such tuples is $\binom{\delta d+k-1}{k-1}\leq(\delta d)^{k}$.





\par\medskip\noindent\textbf{Corollary 5.5.}\par



For each weight $\mu=(\mu_{1},\dots,\mu_{k})$ with $\sum_{i=1}^{k}\mu_{i}=\delta d$ there is an element $H_{\mu}\in\mathcal{U}(\mathfrak{gl}_{k})$ of length at most $(\delta d)^{k}$ which acts on $\mathbb{C}[\mathbb{C}[x_{1},\ldots,x_{k}]_{d}]_{\delta}$ as the projector onto the weight space corresponding to $\mu$.







\par\medskip\noindent\textbf{5.1.2 Isotypic projectors}\par




Let $V$ be a representation of a complex reductive Lie group $G$ with the corresponding Lie algebra $\mathfrak{g}$.
Then $V$ is completely reducible as a representation of $\mathfrak{g}$, with the same isotypic components as for $G$.
As discussed in Section 4.4, they are also isotypic components of $V$ as a representation of $\mathcal{Z}(\mathfrak{g})$, which is a commutative algebra generated by Casimir elements $C_{1},\dots,C_{r}$.





Applying Theorem 5.2, we obtain for every highest weight $\lambda$ a corresponding element $Z_{\lambda}\in\mathcal{U}(\mathfrak{g})$ such that $Z_{\lambda}.v$ is the projection of $v$ onto the isotypic component of corresponding to the irreducible representation with highest weight $\lambda$.
The projectors $Z_{\lambda}$ are obtained as polynomial expressions $Z_{\lambda}=f_{\lambda}(C_{1},\dots,C_{r})$ with $\deg f_{\lambda}$ is at most $p-1$, where $p$ is the diversity of $V$ as a representation of $\mathfrak{g}$.
It follows that if the Casimir elements $C_{1},\dots,C_{r}$ have length at most $\ell$, then the length of $Z_{\lambda}$ is at most $\ell(p-1)$.





For the Lie algebra $\mathfrak{gl}_{k}$ Casimir elements have length at most $k$.
Since every highest weight is a weight, the diversity of $\mathbb{C}[\mathbb{C}[x_{1},\dots,x_{k}]_{d}]_{\delta}$ is bounded by $(d\delta)^{k}$ as in the previous section.





\par\medskip\noindent\textbf{Corollary 5.6.}\par



For each partition $\lambda\mathchoice{\scalebox{1.2}[0.7]{ $\vdash$}_{k}\,}{\scalebox{1.2}[0.7]{ $\vdash$}_{k}\,}{\scalebox{0.7}{\scalebox{1.2}[0.7]{$\vdash$}}_{k}\,}{\scalebox{0.5}{\scalebox{1.2}[0.7]{$\vdash$}}_{k}\,}{\delta d}$ there is an element $Z_{\lambda}\in\mathcal{U}(\mathfrak{gl}_{k})$ of length at most $k(\delta d)^{k}$ which acts on $\mathbb{C}[\mathbb{C}[x_{1},\ldots,x_{k}]_{d}]_{\delta}$ as the projector onto the isotypic component corresponding to $\lambda$.







\par\medskip\noindent\textbf{5.1.3 Highest weight projectors}\par




Since the highest weight subspace of weight $\lambda$ is exactly the weight $\lambda$ subspace of the isotypic component corresponding to the irreducible representation $V_{\lambda}$, the projector onto the highest weight subspace is the composition $H_{\lambda}Z_{\lambda}$ of the projectors constructed in the previous sections.





\par\medskip\noindent\textbf{Corollary 5.7.}\par



For each partition $\lambda\mathchoice{\scalebox{1.2}[0.7]{ $\vdash$}_{k}\,}{\scalebox{1.2}[0.7]{ $\vdash$}_{k}\,}{\scalebox{0.7}{\scalebox{1.2}[0.7]{$\vdash$}}_{k}\,}{\scalebox{0.5}{\scalebox{1.2}[0.7]{$\vdash$}}_{k}\,}{\delta d}$ there is an element $X_{\lambda}=H_{\lambda}Z_{\lambda}\in\mathcal{U}(\mathfrak{gl}_{k})$ of length at most $(k+1)(\delta d)^{k}$ which acts on $\mathbb{C}[\mathbb{C}[x_{1},\ldots,x_{k}]_{d}]_{\delta}$ as the projector onto the highest weight space of weight $\lambda$.







\par\medskip\noindent\textbf{5.1.4 GZ-isotypic projectors}\par




Let $V$ be a representation of a complex reductive Lie group $G$ with the corresponding Lie algebra $\mathfrak{g}$ admitting a Gelfand–Tsetlin chain $\mathfrak{g}_{1}\subset\mathfrak{g}_{2}\subset\dots\subset\mathfrak{g}_{k}=\mathfrak{g}$.
As discussed in Section 4.5, under the action of the Gelfand–Tsetlin algebra $\mathcal{G}\hskip-1.00006pt\mathcal{Z}(\mathfrak{g})$ the representation $V$ decomposes into GZ-isotypic components indexed by Gelfand–Tsetlin patterns, and Theorem 5.2 can be used to construct for every Gelfand–Tsetlin pattern $T$ an element $Y_{T}\in\mathcal{U}(\mathfrak{g})$ such that $Y_{T}.x$ is the projection of $x$ onto the $T$-isotypic subspace of $V$.
The length of $Y_{T}$ is bounded in terms of the lengths of Casimir elements of $\mathfrak{g}_{\ell}$ and the number of Gelfand–Tsetlin patterns appearing in the decomposition of $V$.





Consider the case of $\mathfrak{gl}_{k}$ acting on $\mathbb{C}[\mathbb{C}[x_{1},\dots,x_{k}]_{d}]_{\delta}$.
For the chain $\mathfrak{gl}_{1}\subset\mathfrak{gl}_{2}\subset\dots\subset\mathfrak{gl}_{k}$, Casimir elements have length at most $k$.
Every irreducible subrepresentation of $\mathbb{C}[\mathbb{C}[x_{1},\dots,x_{k}]_{d}]_{\delta}$ corresponds to a partition $\lambda\mathchoice{\scalebox{1.2}[0.7]{ $\vdash$}_{k}\,}{\scalebox{1.2}[0.7]{ $\vdash$}_{k}\,}{\scalebox{0.7}{\scalebox{1.2}[0.7]{$\vdash$}}_{k}\,}{\scalebox{0.5}{\scalebox{1.2}[0.7]{$\vdash$}}_{k}\,}{d\delta}$.
Gelfand–Tsetlin patterns are in bijective correspondence with semistandard Young tableaux with $\delta d$ cells and $k$ rows.
Each row of a tableau is an ordered tuple of integers from $\{1,\dots,k\}$, so there are at most $\binom{d\delta+k-1}{k-1}\leq(d\delta)^{k}$ tuples that can serve as rows of a tableau, and therefore at most $(d\delta)^{k^{2}}$ possible semistandard tableaux.





\par\medskip\noindent\textbf{Corollary 5.8.}\par



For each semistandard tableau $T$ of shape $\lambda\mathchoice{\scalebox{1.2}[0.7]{ $\vdash$}_{k}\,}{\scalebox{1.2}[0.7]{ $\vdash$}_{k}\,}{\scalebox{0.7}{\scalebox{1.2}[0.7]{$\vdash$}}_{k}\,}{\scalebox{0.5}{\scalebox{1.2}[0.7]{$\vdash$}}_{k}\,}{\delta d}$ there is an element $Y_{T}\in\mathcal{U}(\mathfrak{gl}_{k})$ of length at most $k(\delta d)^{k^{2}}$ which acts on $\mathbb{C}[\mathbb{C}[x_{1},\ldots,x_{k}]_{d}]_{\delta}$ as the projector onto the $T$-isotypic space corresponding to $T$.








\par\medskip\noindent\textbf{5.2 Circuit transformations}\par




We now combine the bounds on the lengths of the projectors with the PBW theorem (Theorem 4.9) to prove Theorem 1.1.
Pick any ordered basis $(X_{1},\ldots,X_{k^{2}})$ for $\mathfrak{gl}_{k}$, such as $\{E_{i,j}\}_{i,j}$ with $(i,j)$ ordered lexicographically.





We first introduce some notation.
Set $K\coloneqq k^{2}$.
For $\bm{i}=(i_{1},\ldots,i_{K})\in\mathbb{N}^{K}$ we will use the shorthand notation $X^{\bm{i}}\coloneqq X_{1}^{i_{1}}\cdots X_{K}^{i_{K}}\in\mathcal{U}(\mathfrak{gl}_{k})$.
Let $\bm{i}^{\prime}\leq\bm{i}$ denote that $\bm{i}^{\prime}=(i^{\prime}_{1},\ldots,i^{\prime}_{K})\in\mathbb{N}^{K}$ with $i^{\prime}_{r}\leq i_{r}$ for all $r\in[K]$.
Write $\bm{i}-\bm{i}^{\prime}\coloneqq(i_{1}-i^{\prime}_{1},\ldots,i_{K}-i^{\prime}_{K})$ and $\binom{\bm{i}}{\bm{i}^{\prime}}\coloneqq\binom{i_{1}}{i^{\prime}_{1}}\cdots\binom{i_{K}}{i^{\prime}_{K}}$. Lastly, define $|\bm{i}|\coloneqq\sum_{r}i_{r}$
and $\mathbb{N}^{K}_{\leq L}\coloneqq\{\bm{i}\in\mathbb{N}^{K}\mid|\bm{i}|\leq L\}$.





By ˜4.2, any element $X\in\mathcal{U}(\mathfrak{gl}_{k})$ acts as a derivation on the space of metapolynomials $\mathbb{C}[\mathbb{C}[x_{1},\ldots,x_{k}]_{d}]=\bigoplus_{\delta\geq 0}\mathbb{C}[\mathbb{C}[x_{1},\ldots,x_{k}]_{d}]_{\delta}$.





\par\medskip\noindent\textbf{Lemma 5.9.}\par



Suppose $\Phi,\Gamma\in\mathbb{C}[\mathbb{C}[x_{1},\ldots,x_{k}]_{d}]$ are metapolynomials. Then



\begin{enumerate}


\item 
 


$X^{\bm{i}}.(\Phi+\Gamma)=X^{\bm{i}}.\Phi+X^{\bm{i}}.\Gamma$





\item 
 


$\displaystyle X^{\bm{i}}.(\Phi\Gamma)=\sum_{{\bm{i}}^{\prime}\colon{\bm{i}}^{\prime}\leq{\bm{i}}}\Bigl{(}\begin{array}[]{@{}c@{}}{\bm{i}}\\ {\bm{i}}^{\prime}\end{array}\Bigr{)}\big{(}X^{{\bm{i}}^{\prime}}.\Phi\big{)}\big{(}X^{{\bm{i}}-{\bm{i}}^{\prime}}.\Gamma\big{)}$





\end{enumerate}





\par\medskip\noindent\textbf{Proof.}\par




Item 1 follows directly from linearity.
Since $\mathfrak{gl}_{k}$ acts by derivations (see ˜4.2) we have $X.(fg)=(X.f)g+f(X.g)$.
For $i\in\mathbb{N}$ a straightforward induction argument shows $X^{i}.(fg)=\sum_{i^{\prime}\colon i^{\prime}\leq i}\binom{i}{i^{\prime}}\big{(}X^{i^{\prime}}.f\big{)}\big{(}X^{i-i^{\prime}}.g\big{)}$. Then Item 2 follows.
$\hfill\blacktriangleleft$





\par\medskip\noindent\textbf{Theorem 5.10.}\par



Let $P\in\mathcal{U}(\mathfrak{gl}_{k})$ be an element of length $L$.
Suppose there exists an arithmetic circuit $C$ of size $s$ computing a metapolynomial
$\Delta\colon\mathbb{C}[x_{1},\ldots,x_{k}]_{d}\to\mathbb{C}$.
Then there exists an arithmetic circuit $C^{\prime}$ of size $O\big{(}sL^{2k^{2}}\big{)}$ computing $P.\Delta$.







\par\medskip\noindent\textbf{Proof.}\par




Again set $K\coloneqq k^{2}$.
Construct $C^{\prime}$ from $C$ using the following recursive procedure:



\begin{enumerate}


\item 
 


Add for every input gate $\Theta$
in $C$ the input gates
$X^{\bm{i}}.\Theta$ for ${\bm{i}}\in\mathbb{N}^{K}_{\leq L}$.





\item 
 


Add for every addition gate $\Theta=\Phi+\Gamma$ in $C$ the addition gates computing
$X^{\bm{i}}.\Theta$ for ${\bm{i}}\in\mathbb{N}^{K}_{\leq L}$, using the expression in Lemma 5.9.





\item 
 


Add for every product gate $\Theta=\Phi\Gamma$ in $C$:



\begin{itemize}


\item 
 


product gates computing
$\big{(}X^{{\bm{i}}^{\prime}}.\Phi\big{)}\big{(}X^{{\bm{i}}-{\bm{i}}^{\prime}}.\Gamma\big{)}$
for ${\bm{i}}\in\mathbb{N}^{K}_{\leq L}$ and ${\bm{i}}^{\prime}\leq{\bm{i}}$,





\item 
 


addition gates computing
$X^{\bm{i}}.\Theta$ for ${\bm{i}}\in\mathbb{N}^{K}_{\leq L}$, using the expression in Lemma 5.9.





\end{itemize}




\item 
 


The resulting circuit computes $X^{\bm{i}}.\Delta$ for ${\bm{i}}\in\mathbb{N}^{K}_{\leq L}$.
By the PBW theorem (Theorem 4.9), there exists scalars $\beta_{\bm{i}}\in\mathbb{C}$ such that



\[\begin{aligned}
\displaystyle P=\sum_{\bm{i}\in\mathbb{N}^{K}_{\leq L}}\beta_{\bm{i}}X^{\bm{i}}.
\end{aligned}\]


Add to $C^{\prime}$ the addition gates computing $P.\Delta$.





\end{enumerate}


The largest blow-up occurs in step 3.
Let $A\coloneqq\{({\bm{i}},{\bm{i}}^{\prime})\mid{\bm{i}}\in\mathbb{N}^{K}_{\leq L},{\bm{i}}^{\prime}\leq{\bm{i}}\}$.
We added for every product gate in $C$ exactly $|A|$ product gates computing $\big{(}X^{{\bm{i}}^{\prime}}.f\big{)}\big{(}X^{{\bm{i}}-{\bm{i}}^{\prime}}.g\big{)}$ for $(\bm{i},\bm{i}^{\prime})\in A$,
and $|A|-1$ addition gates computing $X^{\bm{i}}.(fg)=\sum_{(\bm{i},\bm{i}^{\prime})\in A}\binom{{\bm{i}}}{{\bm{i}}^{\prime}}\big{(}X^{{\bm{i}}^{\prime}}.f\big{)}\big{(}X^{{\bm{i}}-{\bm{i}}^{\prime}}.g\big{)}$.





The blow-up of the other steps is bounded above by $|A|$, and the total blow-up is therefore upper bounded by $2|A|-1$.
We now prove that
$|A|=\binom{L+2K}{2K}$.





Let $(\bm{i},\bm{i}^{\prime})\in A$. By definition, $i_{1}+\cdots+i_{K}\leq L$, and $(i_{j}-i_{j}^{\prime})\geq 0$, for all $j\in[K]$, with $i_{j}^{\prime}\geq 0$. Let $a_{j}:=(i_{j}-i_{j}^{\prime})$, and $b_{j}:=i_{j}^{\prime}$. We want to find the number of $(\bm{a},\bm{b})\in\mathbb{N}^{2K}$ such that $\sum_{j=1}^{2K}(a_{j}+b_{j})\leq L$. So the number of pairs $(\bm{a},\bm{b})\in\mathbb{N}^{2K}$ such that $\sum_{j=1}^{2K}(a_{j}+b_{j})=t$ is $\binom{t+2K-1}{2K-1}$. Therefore, by the hockey-stick identity,



\[\begin{aligned}
|A|\;=\;\sum_{t=0}^{L}\binom{t+2K-1}{2K-1}\;=\;\binom{L+2K}{2K}.
\end{aligned}\]




The above equality readily gives us the upper bound of $|A|=O(L^{2K})$, which gives the desired complexity.
$\hfill\blacktriangleleft$






\par\medskip\noindent\textbf{Proof of Theorem 1.1.}\par




Let $\Delta$ be a metapolynomial of format $(d,\delta,n)$ admitting an algebraic circuit of size $s$. For each case of Theorem 1.1, we have a corresponding projector:



\begin{enumerate}


\item 
 


For every weight $\mu$, Corollary 5.5 provides an element $H_{\mu}\in\mathcal{U}(\mathfrak{gl}_{k})$ of length at most $O((\delta d)^{k})$ acting as projection on the weight space of weight $\mu$. Theorem 5.10 guarantees that $H_{\mu}.\Delta$ admits a circuit of size at most $O(s((\delta d)^{k})^{2k^{2}})=O(s(\delta d)^{2k^{3}})$.





\item 
 


For every partition $\lambda$, Corollary 5.6 provides an element $Z_{\mu}\in\mathcal{U}(\mathfrak{gl}_{k})$ of length at most $O(k(\delta d)^{k})$ acting as projection on the isotypic component of type $\lambda$. Theorem 5.10 guarantees that $Z_{\lambda}.\Delta$ admits a circuit of size at most $O(s(k(\delta d)^{k})^{2k^{2}})=O(sk^{2k^{2}}(\delta d)^{2k^{3}})$.





\item 
 


For every partition $\lambda$, Corollary 5.7 provides an element $X_{\lambda}\in\mathcal{U}(\mathfrak{gl}_{k})$ of length at most $O((k+1)(\delta d)^{k})$ acting as projection on the highest weight space of weight $\lambda$. Theorem 5.10 guarantees that $X_{\lambda}.\Delta$ admits a circuit of size at most $O(s((k+1)(\delta d)^{k})^{2k^{2}})=O(s(k+1)^{2k^{2}}(\delta d)^{2k^{3}})$.





\item 
 


For every semistandard tableaux $T$ of shape $\lambda\mathchoice{\scalebox{1.2}[0.7]{ $\vdash$}\,}{\scalebox{1.2}[0.7]{ $\vdash$}\,}{\scalebox{0.7}{\scalebox{1.2}[0.7]{$\vdash$}}\,}{\scalebox{0.5}{\scalebox{1.2}[0.7]{$\vdash$}}\,}d\delta$, Corollary 5.8 provides an element $Y_{\lambda}\nobreak\in\nobreak\mathcal{U}(\mathfrak{gl}_{k})$ of length at most $O(k(\delta d)^{k^{2}})$ acting as projection on $T$-isotypic space. Theorem 5.10 guarantees that $Y_{T}.\Delta$ admits a circuit of size at most $O(s(k(\delta d)^{k^{2}})^{2k^{2}})=O(sk^{2k^{2}}(\delta d)^{2k^{4}})$.





\end{enumerate}


$\hfill\blacktriangleleft$





